% Fuentes: LOGBOOK.md EXP-013, EXP-019..EXP-022, EXP-024, EXP-025, EXP-027; DEC-013.
\section{Técnicas de regularización y resultados}
Cada experimento añade o cambia una sola técnica sobre el modelo anterior y se repite con \res{final-seeds} semillas;
el ruido entre semillas es de unos \res{seed-noise-pp}~pp~\cite{bouthillier}, así que las diferencias menores no se
consideran concluyentes. El orden sigue las preguntas que va dejando cada resultado. Primero se prueba cada técnica
por separado sobre el modelo de la entrega anterior (Cuadro~\ref{tab:phase1}); después, la \emph{data augmentation}
sobre el ganador (Cuadro~\ref{tab:aug}) y las penalizaciones combinadas con ella (Cuadro~\ref{tab:combine});
\emph{label smoothing} sobre el modelo final (Cuadro~\ref{tab:ls}); si Dropout y L2 fallaban por la red pequeña,
variando la capacidad (Sección~\ref{sec:capacity}); y, como la \emph{augmentation} era lo que funcionaba,
desplazamientos para enseñar también la traslación (Sección~\ref{sec:shift}). Como con regularización el entrenamiento es más lento,
el tope de épocas sube de \res{phase0-epochs-cap} a \res{dropout01-epochs-cap}; el modelo de partida no cambia,
porque paraba por \emph{EarlyStopping} mucho antes. El Cuadro~\ref{tab:phases} resume lo que se mantuvo en cada paso.

\begin{table}[h]
  \centering\small
  \safetable{generated/tables/phase_summary}
  \caption{Configuración que se mantiene tras cada paso (validación; media $\pm$ desviación típica).}
  \label{tab:phases}
\end{table}

\begin{table}[h]
  \centering
  \resizebox{\linewidth}{!}{\safetable{generated/tables/phase1_variants}}
  \caption{Cada técnica por separado sobre el modelo de la entrega 1 (\res{final-seeds} semillas por fila; brecha =
  \emph{accuracy} de train $-$ \emph{accuracy} de validación).}
  \label{tab:phase1}
\end{table}

\subsection{Dropout}
En cada paso de entrenamiento, Dropout anula cada unidad oculta con probabilidad $p$ y escala las demás por
$1/(1-p)$; en inferencia se usa la red completa. Equivale a entrenar a la vez un conjunto exponencial de subredes que
comparten pesos y promediarlas al predecir, y evita que las unidades dependan unas de otras
(co-adaptación)~\cite{srivastava,goodfellow}.

\decisionpoint{Si la brecha es sobreajuste, Dropout tras cada capa oculta debería cerrarla subiendo la validación.}
{Al aumentar $p$ la brecha se cierra (de \res{phase0-gap-acc}~pp a \res{dropout05-gap-acc}~pp con
$p$~$=$~\res{dropout05-dropout}), pero porque el \emph{mean recall} de entrenamiento cae
(\res{phase0-train-mean-recall} $\rightarrow$ \res{dropout05-train-mean-recall}), y el de validación cae con él
(Figura~\ref{fig:dropout}).}
{Ninguna intensidad mejora la validación. Con $p$~$=$~\res{dropout01-dropout} el \emph{mean recall} ya queda en
\res{dropout01-mean-recall} (\res{loss-dropout01}~pp), y con $p$~$=$~\res{dropout05-dropout}, en
\res{dropout05-mean-recall} (\res{loss-dropout05}~pp); la clase más pequeña es la primera en perderse: \res{class-min-name} (\res{valid-count-helipad} imágenes de validación) pasa de recall
\res{phase0-recall-helipad} a \res{dropout05-recall-helipad}. Es el comportamiento de un modelo sin capacidad que
sobre: la red de \res{phase0-hidden} ya está cerca del subajuste, y quitarle unidades al azar la empuja a él.}
{Sin Dropout.}

\subsection{Penalización L2}
La penalización L2 añade $\lambda \lVert \mathbf{w} \rVert^2$ a la pérdida, de modo que cada paso encoge los pesos
hacia cero (\emph{weight decay}); equivale a una estimación MAP con una distribución a priori gaussiana sobre los
pesos~\cite{goodfellow}. Con Adam no es exactamente \emph{weight decay}: el gradiente de la penalización se divide,
como el resto, por la raíz de su segundo momento, así que los pesos con gradientes grandes se encogen
menos~\cite{loshchilov}.

\decisionpoint{Penalizar la norma de los pesos limita la complejidad efectiva de la red y debería reducir la brecha.}
{Con $\lambda$~$=$~\res{l2-1e-5-l2} y $\lambda$~$=$~\res{l2-1e-4-l2} nada cambia: la brecha sigue en
\res{l2-1e-5-gap-acc} y \res{l2-1e-4-gap-acc}~pp, y el \emph{mean recall} (\res{l2-1e-5-mean-recall} y
\res{l2-1e-4-mean-recall}) empata con el de partida. Con $\lambda$~$=$~\res{l2-1e-3-l2}, la brecha se cierra
(\res{l2-1e-3-gap-acc}~pp), pero bajando el \emph{mean recall} de entrenamiento a \res{l2-1e-3-train-mean-recall} y el
de validación a \res{l2-1e-3-mean-recall} (Figura~\ref{fig:l2}).}
{La curva de validación de L2 tiene la misma forma que la de Dropout: no hay ningún valor intermedio en que la
validación mejore. La brecha no responde a penalizaciones, lo que apunta a que el límite está en el sesgo del modelo,
una red densa sobre píxeles sin invarianzas espaciales~\cite{goodfellow}, y no en su varianza.}
{Sin L2.}

\begin{figure}[!htbp]
  \centering
  \begin{minipage}[t]{0.485\linewidth}
    \centering
    \safegraphics[width=\linewidth]{figures/dropout_curve.png}{Curva de validación del dropout}
    \caption{\figcaption{dropout_curve}}
    \label{fig:dropout}
  \end{minipage}\hfill
  \begin{minipage}[t]{0.485\linewidth}
    \centering
    \safegraphics[width=\linewidth]{figures/l2_curve.png}{Curva de validación de L2}
    \caption{\figcaption{l2_curve}}
    \label{fig:l2}
  \end{minipage}
\end{figure}

\subsection{BatchNorm}
BatchNorm normaliza la preactivación de cada unidad con la media y la varianza del \emph{batch} y después la escala y
la desplaza con dos parámetros aprendidos ($\gamma$, $\beta$); en inferencia usa medias móviles de esas
estadísticas~\cite{ioffe}. El artículo original lo coloca antes de la no linealidad, y entonces el sesgo de la capa
\texttt{Dense} sobra, porque $\beta$ lo sustituye. La salida no depende de la escala de los pesos, lo que permite
\emph{learning rates} mayores~\cite{ioffe,santurkar}; y como las estadísticas de cada \emph{batch} son ruidosas, añade
algo de regularización~\cite{ioffe}.

\decisionpoint{BatchNorm regulariza y, sobre todo, permite entrenar con una \emph{learning rate} mayor.}
{Con la \emph{learning rate} de partida, BatchNorm antes o después de la activación empata con el modelo sin él
(\res{bn-samelr-mean-recall} y \res{bn-after-mean-recall}; Figura~\ref{fig:bn}). Lejos de regularizar, sube el
\emph{mean recall} de entrenamiento (+\res{gain-bn-train-mr}~pp) y la brecha crece a \res{bn-samelr-gap-acc}~pp: aquí
actúa sobre la optimización, no sobre la generalización.}
{Antes de la activación y con la \emph{learning rate} triplicada (\res{phase1-lr}), el \emph{mean recall} sube
+\res{gain-bn}~pp y el recall mínimo, de \res{phase0-min-recall} a \res{phase1-min-recall}; las tres semillas igualan o
superan a la mejor del modelo de partida, pero la mejora es del orden del ruido. Después de la activación, triplicar
la \emph{learning rate} empeora (\res{bn-after-lr3x-mean-recall}). Su peso real aparece al combinarlo con la
\emph{augmentation}: sin él se pierden \res{loss-no-bn}~pp (Sección~\ref{sec:combine}).}
{BatchNorm antes de la activación, \emph{learning rate} \res{phase1-lr}.}

\subsection{Data augmentation}
La \emph{data augmentation} amplía el conjunto de entrenamiento con transformaciones que conservan la
etiqueta~\cite{goodfellow}. Las imágenes aéreas se toman desde arriba y no tienen una orientación canónica: un avión
girado 90$^\circ$ o reflejado sigue siendo un avión. Las 8 simetrías del cuadrado (grupo diédrico D4: cuatro
rotaciones de 90$^\circ$, con y sin reflejo) conservan la etiqueta y, sobre una rejilla de píxeles, son permutaciones
exactas, sin interpolación. Una red densa no tiene ninguna de estas invarianzas: tiene que aprender cada orientación
por separado, y mostrársela en entrenamiento es la forma de enseñársela. Se comparan volteos, rotaciones de
90$^\circ$, D4 completo y variaciones de brillo y contraste (fotométricas), todas aplicadas al azar en cada
\emph{batch} de entrenamiento, sobre el modelo con BatchNorm y con paciencias de \res{phase2-es-patience} y
\res{phase2-rlr-patience} épocas (Cuadro~\ref{tab:aug}).

\decisionpoint{Las transformaciones geométricas codifican una invarianza real de los datos, así que deberían reducir
la brecha subiendo la validación, no solo bajando el entrenamiento.}
{Con las transformaciones geométricas, el \emph{mean recall} de entrenamiento baja (rotaciones:
\res{no-aug-train-mean-recall} $\rightarrow$ \res{phase2-train-mean-recall}) y el de validación sube
(\res{no-aug-mean-recall} $\rightarrow$ \res{phase2-mean-recall}); la mejor época llega más tarde
(\res{no-aug-best-epoch} $\rightarrow$ \res{phase2-best-epoch} de media), porque hay más variedad que aprender
(Figura~\ref{fig:aug}). Las variaciones de brillo y contraste no reducen la brecha
(\res{color-gap-acc}~pp) y empeoran la validación (\res{loss-color}~pp): la red tiene que aprender a ignorar una
información que, en imágenes de \res{phase2-res}~$\times$~\res{phase2-res} píxeles, es buena parte de la que tiene
(explicación plausible, no verificada).}
{Las rotaciones de 90$^\circ$ suben el \emph{mean recall} +\res{gain-rot90}~pp y la \emph{accuracy}
+\res{gain-rot90-acc}~pp, y reducen la brecha a \res{phase2-gap-acc}~pp. Los volteos (+\res{gain-flip}~pp) y D4
(+\res{gain-d4}~pp) empatan con ellas dentro del ruido. Por clases, ganan sobre todo Cargo plane (\res{no-aug-recall-cargo-plane} $\rightarrow$
\res{phase2-recall-cargo-plane}), Motorboat (\res{no-aug-recall-motorboat} $\rightarrow$
\res{phase2-recall-motorboat}) y Dump truck (\res{no-aug-recall-dump-truck} $\rightarrow$
\res{phase2-recall-dump-truck}); Truck pierde (\res{no-aug-recall-truck} $\rightarrow$ \res{phase2-recall-truck}).}
{Rotaciones de 90$^\circ$ en entrenamiento: la mejor media en las dos métricas y un subconjunto más simple que D4,
aunque la diferencia con D4 y con los volteos no es concluyente.}

\begin{table}[h]
  \centering
  \resizebox{\linewidth}{!}{\safetable{generated/tables/augmentation_variants}}
  \caption{\emph{Data augmentation} sobre el modelo con BatchNorm (\res{final-seeds} semillas por fila; paciencias de
  \res{phase2-es-patience} y \res{phase2-rlr-patience} épocas).}
  \label{tab:aug}
\end{table}

\begin{figure}[!htbp]
  \centering
  \begin{minipage}[t]{0.485\linewidth}
    \centering
    \safegraphics[width=\linewidth]{figures/batchnorm.png}{BatchNorm}
    \caption{\figcaption{batchnorm}}
    \label{fig:bn}
  \end{minipage}\hfill
  \begin{minipage}[t]{0.485\linewidth}
    \centering
    \safegraphics[width=\linewidth]{figures/augmentation.png}{Data augmentation}
    \caption{\figcaption{augmentation}}
    \label{fig:aug}
  \end{minipage}
\end{figure}

\subsection{Combinaciones}\label{sec:combine}
\decisionpoint{Con la \emph{augmentation} cerrando parte de la brecha, una penalización suave podría sumar; y
BatchNorm, que por sí solo aportaba poco, podría pesar más combinado con ella.}
{Sobre BatchNorm + rotaciones (Cuadro~\ref{tab:combine}), Dropout vuelve a bajar entrenamiento y validación a la vez
(\res{loss-c-dropout01} y \res{loss-c-dropout02}~pp con $p$~$=$~\res{c-dropout01-dropout} y
$p$~$=$~\res{c-dropout02-dropout}); una posible causa es que Dropout cambia la varianza de las activaciones entre
entrenamiento e inferencia y desajusta las estadísticas que BatchNorm ha acumulado~\cite{li2019}. L2 deja la
\emph{accuracy} igual y el \emph{mean recall} algo por debajo, dentro del ruido. Es lo esperable: delante de
BatchNorm, la salida no depende de la escala de los pesos, así que L2 no restringe la función que calcula la red;
solo cambia la \emph{learning rate} efectiva~\cite{laarhoven}.}
{Sin BatchNorm (con la \emph{learning rate} de la entrega 1, \res{phase0-lr}, porque la triple se eligió para la
red con BatchNorm), el \emph{mean recall} queda en \res{c-no-bn-mean-recall} (\res{loss-no-bn}~pp), cuando sin
\emph{augmentation} BatchNorm solo aportaba +\res{gain-bn}~pp: con
\emph{augmentation} hay más variedad que ajustar y la optimización más rápida que permite BatchNorm se nota. Quitar la
\emph{augmentation} cuesta \res{gain-rot90}~pp.}
{BatchNorm + rotaciones, sin Dropout ni L2.}

\begin{table}[h]
  \centering
  \resizebox{\linewidth}{!}{\safetable{generated/tables/combine_variants}}
  \caption{Cambios sobre el modelo con BatchNorm y rotaciones de 90$^\circ$ (\res{final-seeds} semillas por fila).}
  \label{tab:combine}
\end{table}

\subsection{Label smoothing}
\emph{Label smoothing} regulariza la salida en lugar de los pesos: sustituye la etiqueta \emph{one-hot} por
$(1-\varepsilon)$ en la clase correcta y $\varepsilon/K$ repartido entre las $K$ clases, de modo que la red ya no gana
nada haciendo crecer sus \emph{logits} sin límite ni dando probabilidad 1 a los ejemplos de
entrenamiento~\cite{szegedy,muller}.

\decisionpoint{Penalizar la sobreconfianza en entrenamiento debería reducir la brecha que queda en el modelo final.}
{Sobre el modelo final (Cuadro~\ref{tab:ls}), con $\varepsilon$~$=$~0,05, 0,1 y 0,2, la brecha baja (de
\res{final-gap-acc}~pp a \res{ls01-gap-acc}~pp con $\varepsilon$~$=$~0,1), pero de nuevo porque baja la
\emph{accuracy} de entrenamiento: la de validación no cambia (\res{ls01-val-acc}, \res{loss-ls01-acc}~pp).}
{El \emph{mean recall} baja con las tres intensidades, con poca dispersión entre semillas (\res{ls01-mean-recall}
con $\varepsilon$~$=$~0,1, \res{loss-ls01}~pp). La pérdida se concentra en pocas clases: \res{class-min-name} pasa de
\res{final-recall-helipad} a \res{ls01-recall-helipad} (solo \res{valid-count-helipad} imágenes de validación, unas
pocas de diferencia) y Storage tank de \res{final-recall-storage-tank} a \res{ls01-recall-storage-tank}, mientras
que Bus y Small car mejoran algo. Repartir probabilidad entre todas las clases resta margen precisamente a las que
tienen pocos ejemplos (explicación plausible, no verificada).}
{Sin \emph{label smoothing}.}

\begin{table}[h]
  \centering
  \resizebox{\linewidth}{!}{\safetable{generated/tables/ls_variants}}
  \caption{\emph{Label smoothing} sobre el modelo final (\res{final-seeds} semillas por fila).}
  \label{tab:ls}
\end{table}

\subsection{Capacidad y regularización}\label{sec:capacity}
Dropout y L2 se juzgaron sobre la red de \res{phase0-hidden}, que la entrega anterior eligió en parte porque las redes
mayores memorizaban: justo lo que la regularización debe corregir. Dropout está pensado para redes con capacidad de
sobra~\cite{srivastava}, así que su veredicto podía depender de la anchura. Para comprobarlo se varían a la vez los
dos factores sobre el modelo final: anchura de las 4 capas ocultas (256, 512 y 1024 unidades, de
\res{final-params} a \res{w1024-params} parámetros) por regularización (ninguna, Dropout 0,2 o 0,5, L2
$10^{-4}$). Es una excepción deliberada a cambiar un solo factor, porque lo que se mide es su interacción.

\decisionpoint{Con más capacidad, Dropout y L2 deberían dejar de empujar al subajuste y ayudar.}
{La interacción existe (Figura~\ref{fig:capreg}): el coste de Dropout 0,2 frente a la misma anchura sin regularizar
pasa de \res{cost-d02-w256}~pp con 256 unidades a \res{cost-d02-w512}~pp con 512 y \res{cost-d02-w1024}~pp con 1024,
y el de Dropout 0,5, de \res{cost-d05-w256} a \res{cost-d05-w1024}~pp. Sin regularizar, la brecha media baja de
\res{final-gap-acc}~pp con 256 unidades a \res{w1024-gap-acc}~pp con 1024, pero no es una diferencia real: la de 256 se
debe sobre todo a una semilla (\res{final-s1-gap-acc}~pp; las otras dos, \res{final-s2-gap-acc} y
\res{final-s3-gap-acc}~pp), y la brecha de cada run depende sobre todo de la época en que se conserva el modelo. Con la
misma validación, una brecha menor solo significa ajustar menos el entrenamiento.}
{Ninguna combinación supera al modelo final. Las redes más anchas sin regularizar empatan con él
(\res{w512-mean-recall} y \res{w1024-mean-recall}, frente a \res{best-mean-recall}), y la mejor celda regularizada,
512 unidades con L2 (\res{w512-l2-mean-recall}~$\pm$~\res{w512-l2-mean-recall-std}), también: empata, con menos
dispersión entre semillas. Dropout reduce su coste al crecer la red, pero no llega a compensarlo en el rango
probado.}
{Se mantiene \res{final-hidden} sin Dropout ni L2: con la validación empatada se elige el modelo más simple, con
varias veces menos parámetros (\res{final-params} frente a \res{w1024-params}). La conclusión de las
secciones anteriores ya no depende de la red pequeña: la regularización explícita tampoco ayuda con más capacidad.}

\begin{figure}[!htbp]
  \centering
  \safegraphics[width=\linewidth]{figures/capacity_reg.png}{Capacidad por regularización}
  \caption{\figcaption{capacity_reg}}
  \label{fig:capreg}
\end{figure}

\subsection{Desplazamientos}\label{sec:shift}
Con la regularización explícita descartada también para redes más anchas, queda la vía que sí funcionó: dar a la red
invarianzas a través de los datos. Las rotaciones le enseñaron la orientación; la invarianza que más le falta a una red
densa es la de traslación, porque cada píxel tiene su propio peso~\cite{goodfellow}. Se añaden desplazamientos
aleatorios de la imagen (rellenando el borde por reflejo) al modelo final.

\decisionpoint{Desplazar la imagen unos píxeles conserva la etiqueta y debería enseñar a la red la invarianza a
traslaciones, como las rotaciones le enseñaron la de orientación.}
{Con desplazamientos de hasta 1 y 2 píxeles (Cuadro~\ref{tab:shift}), el \emph{mean recall} cae
\res{loss-shift1} y \res{loss-shift2}~pp, y también el de entrenamiento (\res{shift1-train-mean-recall} y
\res{shift2-train-mean-recall}): la red vuelve al subajuste. A \res{final-res}~$\times$~\res{final-res}, un píxel es una
fracción grande de la imagen y los recortes están encuadrados sobre el objeto, así que el desplazamiento recorta parte
de él (explicación plausible, no verificada).}
{Para ver si era el tamaño del desplazamiento, se prueba medio píxel. Reducir la imagen a 32 píxeles por área y
promediar después bloques de 2~$\times$~2 equivale a reducirla a 16, así que desplazar un píxel la imagen de 32 antes
de una capa de promedio es desplazar medio píxel la de 16. El control sin desplazamientos queda a
\res{loss-pool-control}~pp del modelo final, del orden del ruido (difieren el redondeo de las imágenes y la
trayectoria del entrenamiento), y frente a ese control el medio píxel cuesta \res{loss-half}~pp. El patrón es
monótono: cuanto mayor el desplazamiento, peor.}
{Sin desplazamientos. A esta resolución, la invarianza a traslaciones no se le puede dar a una red densa a través de
los datos; es justo lo que una red convolucional incorpora en su arquitectura.}

\begin{table}[h]
  \centering
  \resizebox{\linewidth}{!}{\safetable{generated/tables/shift_variants}}
  \caption{Desplazamientos aleatorios sobre el modelo final (\res{final-seeds} semillas por fila). Las dos últimas
  filas usan la entrada de 32 píxeles con un promedio de 2~$\times$~2 delante de la red.}
  \label{tab:shift}
\end{table}
